\documentclass[11pt,letterpaper]{article}
\usepackage[T1]{fontenc}
\usepackage{lmodern,microtype}
\usepackage[margin=1in,headheight=15pt]{geometry}
\usepackage{amsmath,amssymb,amsthm,mathtools}
\usepackage{booktabs,longtable,array}
\usepackage{xcolor,listings,xurl}
\usepackage{enumitem,fancyhdr,needspace,etoolbox}
\usepackage[hidelinks]{hyperref}
\hypersetup{pdftitle={RadicalRoots: Exact radicals from algebraic roots},pdfauthor={OpenAI, prepared for Vladimir Reshetnikov}}
\pagestyle{fancy}\fancyhf{}
\fancyhead[L]{\small RadicalRoots 0.1.0}
\fancyhead[R]{\small Exact radicals from algebraic roots}
\fancyfoot[C]{\thepage}
\emergencystretch=2em
\setlength{\parskip}{0.4em}\setlength{\parindent}{0pt}
\setlist{nosep,leftmargin=1.6em}
\lstset{basicstyle=\ttfamily\small,breaklines=true,columns=fullflexible,
 keepspaces=true,showstringspaces=false,frame=single,framerule=0.3pt,
 aboveskip=0.8em,belowskip=0.8em}
\BeforeBeginEnvironment{lstlisting}{\Needspace{9\baselineskip}}
\newcommand{\Q}{\mathbb Q}\newcommand{\C}{\mathbb C}
\newcommand{\Gal}{\operatorname{Gal}}\newcommand{\rad}{\operatorname{rad}}
\newcommand{\Fix}{\operatorname{Fix}}\newcommand{\id}{\operatorname{id}}
\newcommand{\code}[1]{\texttt{#1}}
\newtheorem{theorem}{Theorem}[section]
\newtheorem{lemma}[theorem]{Lemma}
\newtheorem{proposition}[theorem]{Proposition}
\newtheorem{corollary}[theorem]{Corollary}
\theoremstyle{definition}\newtheorem{definition}[theorem]{Definition}
\theoremstyle{remark}\newtheorem{remark}[theorem]{Remark}
\begin{document}
\begin{titlepage}
\vspace*{1.0cm}
{\Large Computational algebra\par}
\vspace{0.8cm}
{\Huge\bfseries RadicalRoots\par}
\vspace{0.45cm}
{\LARGE Exact radicals from algebraic roots\par}
\vspace{0.75cm}
{\large Structural reductions, constructive Galois theory,\par
and branch-aware reference implementations\par}
\vspace{1cm}
Prepared for Vladimir Reshetnikov\par
8 September 2026\par
Reference implementation 0.1.0
\vspace{1.0cm}
\begin{abstract}
An algebraic \code{Root} object specifies both a polynomial equation and a
particular complex root. Producing radicals therefore requires more than
finding a solvable equation: every radical branch must be made consistent
with that embedding. This article develops a two-layer solution. A fast layer
recognizes polynomial composition, powers, generalized reciprocal symmetry
and Dickson polynomials. A general layer constructs an embedded splitting
field, adjoins a carefully chosen cyclotomic field, computes actual
field automorphisms and applies finite Fourier resolvents along a
prime-index normal series. It produces explicit principal-power expressions
with exactly selected root-of-unity corrections.

The unbounded mathematical construction covers every algebraic number over
\(\Q\) that is expressible by radicals. The accompanying software implements
that construction using computer-algebra primitives; it is a reference
implementation, not an unconditional performance or backend-completeness
claim. The Python implementation has executable regression results. The
Wolfram Language implementation and its native tests are supplied but were
not executed because the connected Wolfram service returned HTTP 404.
\end{abstract}
\vfill
\small
Deliverables: Wolfram Language package; Python/SymPy implementation;
regression tests and recorded results; source-notebook inspection notes;
this article in editable \LaTeX\ and compiled PDF.
\end{titlepage}
\begingroup
\small
\setlength{\parskip}{0pt}
\tableofcontents
\endgroup
\newpage

\section{Problem, contract, and scope}

The motivating question asks why \code{ToRadicals} does not convert certain
explicit algebraic numbers despite known radical formulas \cite{question}.
Wolfram's documentation explicitly notes that some higher-degree solvable
polynomials are not converted by that function \cite{toradicals}.
Consequently, an unchanged \code{Root} is not a proof of nonsolvability.

The input considered here is an exact algebraic number \(\alpha\in\overline\Q\)
with its specified embedding in \(\C\). The ordinary Wolfram example is
\code{Root[f\&,k]} for a rational-coefficient polynomial. The Python front end
accepts a rational polynomial and a zero-based SymPy root index. It first
selects the actual root and then works with its \emph{minimal polynomial}.
This distinction is important for a reducible input: a rational root of a
polynomial having an additional nonsolvable factor is still rational.

\begin{definition}[Output language]
A radical expression is built from rational constants and \(i\), by addition,
multiplication and rational powers, with nonzero denominators. Powers use
principal complex values. Equivalently, division is a negative integer power.
No \code{Root}, \code{CRootOf}, \code{AlgebraicNumber}, symbolic parameter,
inexact floating-point constant, trigonometric function or exponential
function is permitted in a successful output.
\end{definition}

The syntax is deliberately stricter than ``a short expression involving
algebraic special-function values.'' For instance, \(\cos(\pi/7)\) is algebraic
and solvable, but is not an output of the specified grammar. Conversely,
\[
             \zeta_m=(-1)^{2/m}
\]
\emph{is} a permitted radical expression: with the principal logarithm of
\(-1\), its value is \(e^{2\pi i/m}\). The exponential is explanatory notation,
not an output head. The cases \(m=1,2\) give \(1,-1\).

A successful public result must satisfy both the grammar test and the
selected-root equality test. Wolfram uses \code{RootReduce[e-alpha]===0}; the
Python implementation compares exact minimal polynomials and then performs
a separation-bound root-identity test. These are CAS-level checks, not
proof-assistant certificates. Internal \code{Root} and number-field objects
are allowed while constructing and verifying the answer; they must not
remain in the returned expression.

The package does not solve parameter-dependent equations, decide whether an
arbitrary special-function expression is algebraic, minimize radical depth,
or guarantee the shortest printed expression. These are different problems.
The default finite resource limits also mean that a solvable input can yield
a resource-limit report rather than radicals.

\section{What the notebook contributes}

The supplied \code{FindExtension.nb} is an extensive exploratory session.
Inspection of the archive, without evaluating it, located 220 Input cells.
The accompanying diagnostic transcription records those cells and the
original notebook line numbers. It is not presented as an executable export:
some special boxes and quantifier notation remain.

The essential strategy of \code{findExtension} is to divide the minimal
polynomial by a prospective quadratic
\(X^2+pX+q\), eliminate the condition that the remainder vanish, and select
a comparatively simple algebraic value of a coefficient. The cubic variant
\code{findExtension3} applies the same idea to
\(X^3+pX^2+qX+r\). This is a useful way of looking for an extension over which
the original polynomial has a small factor. The subsequent steps factor
over that extension and select the factor containing the desired root.

Additional experiments project roots onto real or imaginary parts, search
subset sums of conjugates, recognize candidate minimal polynomials from
high-precision values, and use manually selected radical or cyclotomic
extensions. Those ideas motivate the fast layer and the extension-hint
interface. They do not, by themselves, constitute an exhaustive algorithm.
A solvable group need not expose the desired tower through quadratic or
cubic factors alone.

Several engineering changes are required to turn this exploration into a
reusable solver. The notebook includes \code{While[True]} searches, selection
of the first candidate without a general empty-result contract, references
to private simplification symbols and stateful constructs such as previous
outputs. Its initial complexity function counts \code{Sin} and \code{Cos}
in a depth measure; a small cost therefore does not imply a radical-only
answer. In two early cells the literal private name
\code{Simplify`SimpifyCount} differs from the later
\code{Simplify`SimplifyCount}. The replacement relies on neither spelling.

Numerical recognition is useful for proposing an algebraic identity, but a
400-digit match is not the same as establishing that identity. Similarly,
\code{PossibleZeroQ} is not the explicit branch certificate chosen here.
The replacement separates proposal generation, exact root selection,
radical-only syntax, resource limits and mathematical nonsolvability.
The source-review directory contains the detailed mapping and archive hashes.

\section{The mathematical existence criterion}

Let \(f\in\Q[X]\) be the monic irreducible minimal polynomial of \(\alpha\),
let \(L\subset\C\) be its splitting field, and put \(G=\Gal(L/\Q)\).
Galois's solvability theorem states that a characteristic-zero polynomial
is solvable by radicals precisely when its Galois group is solvable
\cite[Theorem 5.34]{milne}.

\begin{proposition}[One specified root is enough]
The algebraic number \(\alpha\) is expressible by radicals over \(\Q\) if and
only if \(G\) is solvable.
\end{proposition}
\begin{proof}
If the whole polynomial is solvable by radicals, so is its selected root.
Conversely, suppose \(\alpha\) belongs to a radical tower. Each embedding of
\(\Q(\alpha)\) into an algebraic closure extends to the algebraic fields
containing that tower. The images of its radical generators remain roots of
the corresponding power equations. Thus every conjugate of \(\alpha\) also
belongs to a radical extension. Taking the compositum of the finitely many
such towers still gives a radical extension containing every root of \(f\).
Apply the polynomial solvability theorem.
\end{proof}

This statement requires the \emph{minimal} polynomial. It does not say that
\(\Q(\alpha)/\Q\) is normal, nor that the necessary radical tower must be
contained in \(\Q(\alpha)\). For example, a cubic may require auxiliary
complex radicals even when the selected root is real. Searching only inside
the degree-\(\deg f\) root field is therefore not a complete approach.

For a finite group, solvability can be decided by repeatedly taking
commutator subgroups. The sequence either reaches the identity or stabilizes
at a nontrivial perfect subgroup. The constructive algorithm below refines
this into a normal series with prime cyclic quotients; it needs the actual
field action, not merely the abstract isomorphism type of the group.

\section{Fast structural reductions}
\label{sec:fast}

The fast layer operates on equations, generates candidates, and then checks
which candidate equals the selected input. It does not infer a branch from
a numerical ordering alone. Affine depression is performed first when
needed: for a monic degree-\(n\) polynomial with next coefficient \(a_{n-1}\),
substitute \(X=Y-a_{n-1}/n\).

\subsection{Low degree, powers, and composition}

Degrees one and two use the standard formulas. For a depressed cubic
\(Y^3+PY+Q=0\), set
\[
 \Delta=\frac{Q^2}{4}+\frac{P^3}{27},\qquad
 t=-\frac Q2+\sqrt\Delta,\qquad u=t^{1/3}.
\]
Choose the other sign of \(\sqrt\Delta\) when the first \(t\) is zero, unless
\(P=Q=0\), which is handled separately. The complete coherent candidate list
is
\[
        Y_k=\zeta_3^k u-\frac{P}{3\zeta_3^k u},
        \qquad k=0,1,2.
\]
Indeed, the two summands have product \(-P/3\), and their cubes sum to \(-Q\).
This product constraint is what independently chosen cube roots can violate.
Quartic candidates come from the CAS's explicit quartic solver and are
subsequently checked.

If every nonconstant exponent is divisible by \(r>1\), then
\(f(X)=g(X^r)\). Solve \(g(Y)=0\) recursively and generate
\(X=\zeta_r^kY^{1/r}\), for all \(0\leq k<r\).
The zero case is harmless. Likewise, a rational polynomial decomposition
\[
                 f=f_1\circ f_2\circ\cdots\circ f_s
\]
can be solved from the outside inward. Intermediate coefficients may then
be radicals. The package does not claim to discover every decomposition
over every algebraic coefficient field: rational decomposition is a fast
path, not the completeness mechanism.

\subsection{Dickson polynomials and coherent inverse pairs}

Define
\[
 D_0(X,a)=2,\qquad D_1(X,a)=X,\qquad
 D_n(X,a)=X D_{n-1}(X,a)-aD_{n-2}(X,a).
\]
Direct induction gives the key identity
\begin{equation}
       D_n(u+a/u,a)=u^n+(a/u)^n.                 \label{eq:dickson}
\end{equation}
For \(a\ne0\), an equation \(D_n(X,a)+c=0\) is therefore solved by
\begin{equation}
 \begin{split}
 t&=\frac{-c+\sqrt{c^2-4a^n}}2,\qquad u=t^{1/n},\\
 X_k&=\zeta_n^k u+\frac{a}{\zeta_n^k u},\qquad 0\leq k<n.
 \end{split}                                                 \label{eq:dicksonroots}
\end{equation}
Here \(t\) satisfies \(t^2+ct+a^n=0\), so the two \(n\)-th powers in
\eqref{eq:dickson} sum to \(-c\). Because \(a\ne0\), neither quadratic root
is zero. The code nevertheless includes a zero guard. The case \(a=0\)
reduces to a binomial.

The identity is not a heuristic recognition rule. After depression, the
coefficient of \(X^{n-2}\) suggests
\(a=-[X^{n-2}]f/n\); the implementation then checks the \emph{entire}
polynomial identity \(f-D_n=c\). The same pair \(u,a/u\) is used throughout.
Taking two unrelated principal \(n\)-th roots need not preserve their
required product \(a\).

\subsection{Generalized reciprocal symmetry}

Let \(f(X)=\sum_{k=0}^{2m}c_kX^k\) be monic. Suppose \(a\in\Q^\times\)
satisfies
\begin{equation}
 c_{m-j}=a^j c_{m+j}\quad(1\leq j\leq m).       \label{eq:recipcondition}
\end{equation}
Necessarily \(c_0=a^m\). Grouping the positive and negative powers yields
\[
 \frac{f(X)}{X^m}=c_m+\sum_{j=1}^m c_{m+j}
                     \left(X^j+(a/X)^j\right).
\]
By \eqref{eq:dickson}, this equals \(g(X+a/X)\), where
\[
              g(Y)=c_m+\sum_{j=1}^m c_{m+j}D_j(Y,a).
\]
Thus a degree-\(2m\) problem reduces to a degree-\(m\) equation and quadratics:
\begin{equation}
              X=\frac{Y\pm\sqrt{Y^2-4a}}2.                   \label{eq:inverse-recip}
\end{equation}
The implementation enumerates the rational solutions of \(a^m=c_0\), then
checks every condition in \eqref{eq:recipcondition}. This includes ordinary
reciprocity \(a=1\) and anti-reciprocal substitutions \(a=-1\).

\subsection{Explicit extension hints}

The Wolfram interface also accepts a list of already explicit radical
constants as \code{"ExtensionHints"}. It factors over their compositum,
selects smaller-degree factors by exact vanishing at the original root, and
passes them to the fast solver. This retains the useful part of the notebook
workflow without confusing a guessed extension with a proof of success.
The Python front end does not currently implement this hint interface;
its supplied-field interface serves a different, lower-level purpose.

\section{The two motivating examples}
\label{sec:examples}

\subsection{The sextic: a cubic hidden by \texorpdfstring{\(X-1/X\)}{X-1/X}}

Consider
\[
             f(X)=X^6+X^4-X^3-X^2-1.
\]
The exact identity
\begin{equation}
 f(X)=X^3\left((X-X^{-1})^3+4(X-X^{-1})-1\right)             \label{eq:sextic}
\end{equation}
exhibits generalized reciprocity with \(a=-1\). Put
\begin{equation}
 t=\left(\frac12+\frac{\sqrt{849}}{18}\right)^{1/3},\qquad
 y=t-\frac4{3t},\qquad
 \alpha=\frac{y+\sqrt{y^2+4}}2.                              \label{eq:sexticrad}
\end{equation}
All roots in this displayed formula are positive real principal roots.
The cubic is \(y^3+4y-1=0\): its Cardano discriminant is
\(1/4+64/27=283/108\), whose square root is \(\sqrt{849}/18\).

The cubic derivative \(3y^2+4\) is positive on \(\mathbb R\), so it has exactly
one real root. For that root, \(X^2-yX-1=0\) has one negative and one positive
solution. Nonreal \(y\) cannot come from a real nonzero \(X\). Hence \(f\) has
exactly two real roots, and \(\alpha\) is the positive one: the requested
Wolfram \code{Root[-1-\#\^{}2-\#\^{}3+\#\^{}4+\#\^{}6\&,2]}.
This is a structural explanation of the conversion, not a lucky extension
search. The three cubic candidates, each lifted by both signs in
\eqref{eq:inverse-recip}, yield all six conjugates.

\subsection{The quintic: an inverse Dickson value}

The second polynomial is
\[
           5X^5-25X^3+25X+6=5\left(D_5(X,1)+\frac65\right).
\]
Set
\begin{equation}
              u=\left(\frac{-3+4i}{5}\right)^{1/5},\qquad
              \beta=u+\frac1u.                              \label{eq:quinticrad}
\end{equation}
Here \(u\) is the principal fifth root. Its fifth power has modulus one, so
\(u\) also has modulus one. The inverse-pair identity immediately gives
\(D_5(\beta,1)=-6/5\).

To identify the selected branch analytically, write
\((-3+4i)/5=e^{i\theta}\) with \(\pi/2<\theta<\pi\). Then the five candidates
are \(2\cos((\theta+2\pi k)/5)\), \(0\leq k<5\), and are distinct and real.
The angle \(\theta/5\) is closest to a multiple of \(2\pi\), so \(k=0\)
gives the largest root. Thus \eqref{eq:quinticrad} equals the requested
Wolfram root with index 5. The trigonometric notation is used only in this
proof of ordering; it is absent from the returned radicals.

The Python suite tests every conjugate of both polynomials, not just these
two branches. The sextic's positive real root is Python index 1, and the
quintic's largest real root is index 4. These particular correspondences
must not be generalized to arbitrary complex-root indices across systems.

\section{A complete constructive route through a splitting field}
\label{sec:general}

The remaining construction uses only exact algebraic-number operations and
finite group computations. Its purpose is completeness in principle, not
competition with specialized radical formulas on their natural inputs.

\subsection{Build the normal closure with its embedding}

Obtain every exact root of the irreducible polynomial \(f\), and construct
\[
                 L=\Q(\alpha_1,\ldots,\alpha_n)\subset\C.
\]
A primitive element \(\theta_L\) represents this field as a rational
polynomial quotient. Because all roots of \(f\) are included, the field is
normal; characteristic zero gives separability. Its degree
\(d=[L:\Q]\) is the order of its Galois group.

The Wolfram implementation uses \code{ToNumberField[values,All]}: the
\code{All} is essential because that form is documented to use the smallest
common field extension, whereas the automatic form need not
\cite{tonumberfield}. The Python implementation performs successive exact
primitive-element adjunctions. For degrees at most four, it may use explicit
radical roots \emph{internally} as field generators to avoid a costly
\code{CRootOf} adjunction. That optimization does not replace Fourier descent
when the general method is requested.

\subsection{Adjoin roots of unity before taking resolvents}

Define the squarefree integer
\begin{equation}
        m=\rad(d)=\prod_{p\mid d}p,
        \qquad C=\Q(\zeta_m),\qquad M=LC.                    \label{eq:compositum}
\end{equation}
Use \(m=1\) if \(d=1\). Let
\[
                  H=\Gal(M/C).
\]
Every prime divisor of \(|H|\) divides \(d\), hence divides \(m\). Consequently
\(C\) contains every primitive prime-order root of unity needed for a
prime-quotient normal series of \(H\).

\begin{lemma}[The cyclotomic base does not hide nonsolvability]
The group \(H\) is solvable if and only if \(G=\Gal(L/\Q)\) is solvable.
\end{lemma}
\begin{proof}
Restriction identifies \(H\) with
\(\Gal(L/(L\cap C))\). The intersection \(L\cap C\) is Galois over \(\Q\),
and its Galois group is abelian because it is a subquotient of the
cyclotomic Galois group. Thus \(H\) identifies with a normal subgroup of
\(G\) with abelian quotient. Subgroups of solvable groups are solvable;
conversely, an extension of a solvable group by an abelian group is solvable.
\end{proof}

A common error is to compute automorphisms of \(L/\Q\), adjoin a formal
\(\zeta_p\), and then assume the old automorphisms fix it. Intersections with
the cyclotomic field make that assumption invalid. We construct the actual
compositum and explicitly retain only automorphisms fixing its chosen
\(\zeta_m\). No linear-disjointness assumption is made.

\subsection{Recover actual automorphisms from a primitive element}

Let \(M=\Q(\theta)\), with degree \(D\) and minimal polynomial \(F\).
Every element has a unique coordinate representation
\[
            a=A(\theta),\qquad A\in\Q[T],\quad \deg A<D.
\]
Since \(M/\Q\) is Galois, \(F\) has \(D\) distinct roots \(\theta_j\) in
\(M\). Each gives an automorphism
\begin{equation}
             \sigma_j(A(\theta))=A(\theta_j).                \label{eq:actualaction}
\end{equation}
Composition is determined by the image of \(\theta\): find the unique index
\(k\) for which \(\sigma_i(\theta_j)=\theta_k\). This yields a multiplication
table on actual embedded field maps. Select
\[
              H=\{\sigma_j:\sigma_j(\zeta_m)=\zeta_m\}.
\]
The dimension check
\begin{equation}
                      |H|\varphi(m)=D                       \label{eq:dimension}
\end{equation}
is an inexpensive but important invariant.

Wolfram enumerates the conjugates of \(\theta\) as exact roots and obtains
coordinates using \code{ToNumberField}. Python factors \(F\) over \(M\);
exactly \(D\) distinct linear factors certify normality and provide all
images. Its field elements are rational coefficient arrays reduced modulo
\(F\), and \eqref{eq:actualaction} is evaluated by Horner's rule.

An abstract permutation group returned by a Galois-group routine is not
silently assigned to a numerically ordered root list. The optional SymPy
small-degree group computation is used only to decide solvability, never
to attach an unverified action to root labels.

\section{Constructing a prime cyclic normal series}

We need
\begin{equation}
 H=H_0\triangleright H_1\triangleright\cdots
     \triangleright H_r=1,
 \qquad H_i\triangleleft H_{i-1},\quad
 H_{i-1}/H_i\simeq C_{p_i},                                  \label{eq:chain}
\end{equation}
where each \(p_i\) is prime. This can be found without enumerating the entire
subgroup lattice.

For a current nontrivial subgroup \(K\), generate its derived subgroup
\(N=[K,K]\). If \(N=K\), this nontrivial perfect subgroup proves that the
original group is not solvable. Otherwise enlarge \(N\) greedily by adjoining
an element of \(K\) whenever the generated subgroup remains proper.
Every subgroup encountered contains \([K,K]\), so its quotient image in
\(K/[K,K]\) is a subgroup of an abelian group; it is therefore normal in
\(K\). When no further proper enlargement is possible, \(N\) is maximal
proper and normal. The quotient \(K/N\) is a simple finite abelian group,
hence cyclic of prime order. Repeat within \(N\).

A single pass over the elements of \(K\) suffices for the greedy enlargement.
Once adjoining an element would generate all of \(K\), making \(N\) larger
cannot make that same adjunction proper later. The implementation independently
checks normality, primality of the index, and that the chosen coset generator
\(\sigma\) satisfies \(\sigma^{p}\in N\).

Under Galois correspondence, \eqref{eq:chain} gives
\[
       C=M^{H_0}\subset M^{H_1}\subset\cdots\subset M^{H_r}=M,
\]
with cyclic prime-degree steps. The next section translates these steps
into explicit radicals without choosing a potentially vanishing special
resolvent as a generator.

\section{Fourier descent, zero resolvents, and exact phases}
\label{sec:fourier}

\subsection{The local identity}

Fix one step \(N\triangleleft K\) with \(K/N\) cyclic of prime order \(p\).
Choose a representative \(\sigma\) generating the quotient. Let
\(a\in M^N\), and let \(\zeta=\zeta_p\in C\). Define
\begin{equation}
       R_j=\sum_{k=0}^{p-1}\zeta^{-jk}\sigma^k(a),
       \qquad 0\leq j<p.                                    \label{eq:resolvent}
\end{equation}

\begin{lemma}[Fourier descent]
Each \(R_j\) lies in \(M^N\), \(R_0\in M^K\), and \(R_j^p\in M^K\).
Moreover,
\begin{equation}
       \sigma(R_j)=\zeta^jR_j,
       \qquad a=\frac1p\sum_{j=0}^{p-1}R_j.                  \label{eq:Fourierinverse}
\end{equation}
\end{lemma}
\begin{proof}
Normality of \(N\) shows that every \(\sigma^k(a)\) is fixed by \(N\),
and \(N\) fixes \(\zeta\). Thus the resolvents are fixed by \(N\).
The relation \(\sigma^p\in N\) gives \(\sigma^p(a)=a\).
Shifting the summation index in \eqref{eq:resolvent} now gives the displayed
eigenvalue equation. It makes \(R_0\) and \(R_j^p\) fixed by \(\sigma\),
hence by \(K=\langle N,\sigma\rangle\). Finally,
\(\sum_{j=0}^{p-1}\zeta^{-jk}\) is \(p\) for \(k=0\) and zero otherwise.
Summing \eqref{eq:resolvent} over \(j\) proves the inverse identity.
\end{proof}

The identity uses \emph{all} Fourier components. Some may be zero; they are
simply omitted. There is no search for a nonzero eigenvector and no division
by a resolvent. If \(\sigma(a)=a\), then \(a\) already belongs to the parent
fixed field and the entire step can be skipped.

\subsection{Principal-root selection}

Suppose \(R_j\ne0\), and let \(b=R_j^p\in M^K\). Recursively express \(b\)
by radicals as \(B\). The principal \(p\)-th root \(v=b^{1/p}\) belongs to
\(M\): all roots of \(T^p-b\) are \(\zeta^kR_j\), and all lie in that field.
Therefore there is a unique \(\ell\in\{0,\ldots,p-1\}\) such that
\begin{equation}
                  R_j=\zeta^\ell v.                         \label{eq:phase}
\end{equation}
Test this finite list using exact embedded algebraic equality, and return
\(\zeta_p^\ell B^{1/p}\). Because \(B\) and \(b\) are the same complex number,
their principal powers agree. This is why the selected embedding must be
preserved during recursion.

The phase is computed from the compact internal field value \(b\), not from
a repeatedly expanded radical output. The selected phase is recorded in the
report. Neither a small imaginary part nor a nearest approximate root is
accepted as a certificate. In particular, a negative real cube root generally
requires a nonzero phase relative to the principal complex cube root.

\subsection{The bottom field}

At recursion level zero, \(a\in C\). Write uniquely
\begin{equation}
              a=\sum_{k=0}^{\varphi(m)-1}q_k\zeta_m^k,
              \qquad q_k\in\Q.                              \label{eq:bottom}
\end{equation}
Replace \(\zeta_m\) by \((-1)^{2/m}\). This is already a radical expression,
so recursion ends. The Python implementation solves the coordinate system
for this cyclotomic basis inside \(M\) and checks the full reconstructed
vector, not only a square subsystem. Wolfram uses \code{ToNumberField} with
the actual cyclotomic generator; a rational bottom value is returned directly.

\subsection{Pseudocode}

\begin{lstlisting}
Encode(a, level):
    enforce the node budget; use memoization
    if a = 0: return 0
    if level = 0: return cyclotomic-coordinate expression for a
    (K, N, sigma, p) = chain[level]
    assert a is fixed by N
    if sigma(a) = a: return Encode(a, level - 1)
    compute all R[j] from the p-term orbit of a
    assert R[0] is fixed by K
    terms = [Encode(R[0], level - 1)]
    for j = 1, ..., p-1:
        if R[j] = 0: append 0; continue
        b = R[j]^p
        assert b is fixed by K
        B = Encode(b, level - 1)
        v = principal_pth_root(b), embedded exactly in M
        find ell with zeta_p^ell * v = R[j], exactly
        append zeta_p^ell * principal_pth_root(B)
    return sum(terms)/p
\end{lstlisting}

\section{Correctness and completeness of the construction}

\begin{theorem}[Constructive sufficiency]
Assume the exact field operations, primitive-element computations,
factorizations and embedded algebraic comparisons described above are
available. If \(\alpha\) is solvable by radicals, the unbounded construction
returns a radical expression equal to the specified \(\alpha\).
\end{theorem}
\begin{proof}
The splitting field and cyclotomic compositum are finite, so their exact
construction terminates under the stated primitives. The relative group
\(H\) is solvable by the lemma in Section~\ref{sec:general}. Thus the finite
group procedure reaches the identity and produces \eqref{eq:chain}.

Prove by induction on the recursion level that \code{Encode(a,level)} is a
permitted radical expression equal to \(a\), whenever \(a\) lies in the
corresponding fixed field. At level zero this is \eqref{eq:bottom}.
At a positive level, the Fourier lemma puts every recursively needed value
in the parent field. The induction hypothesis supplies its radical expression;
\eqref{eq:phase} reconstructs each nonzero resolvent with the correct complex
value. The inverse Fourier identity reconstructs \(a\). Zero resolvents and
skipped fixed-field steps preserve the same conclusion. Every recursive call
reduces the level, so the process is finite. At the top level the fixed
subgroup is trivial and the input \(\alpha\in M\) is allowed.
\end{proof}

\begin{theorem}[Decision in the unbounded mathematical model]
For every exact algebraic \(\alpha\) over \(\Q\), the construction either
returns correct radicals or proves that no such radicals exist, provided
the stated exact primitives terminate.
\end{theorem}
\begin{proof}
The solvable case is the previous theorem. Otherwise the relative group is
nonsolvable. The normal-series procedure encounters a nontrivial perfect
subgroup in finitely many strict subgroup reductions. The cyclotomic-base
lemma then implies that the minimal polynomial's Galois group is nonsolvable.
The existence criterion rules out radicals for its specified root.
\end{proof}

\begin{remark}[A theorem is not a blanket software warranty]
The reference implementations invoke large CAS primitives, do not implement
a formally verified algebraic-number kernel, and impose finite default
budgets. A backend exception remains a backend exception even with the
budgets removed. The completeness theorem describes the algorithm with
terminating exact primitives; it is not a claim that SymPy or a particular
Wolfram release successfully executes every such primitive for every input.
\end{remark}

\subsection{The supplied-model caveat}

Python also accepts a separately constructed, validated Galois field model.
It checks normality, cyclotomic containment, actual automorphism actions and
\eqref{eq:dimension}. However, that field can be larger than the target's
normal closure. A nonsolvable supplied model \emph{does not} prove the target
nonsolvable: even a rational number belongs to nonsolvable extensions.
The low-level supplied-model routine therefore raises a backend/model error
in this situation, not a target-level \code{NotSolvable} conclusion. The
automatic construction has the stronger minimal-splitting-field provenance
needed for that conclusion.

\section{Certification details and software architecture}

\subsection{Keep field values and radical syntax separate}

Internal values are canonical algebraic numbers or rational vectors modulo
a primitive polynomial. Output values are ordinary radical expressions.
Each recursion node maintains both views. Expanding the radical expression
is unnecessary for applying automorphisms; operations on the internal value
are smaller and make fixed-field checks direct. Memoization keys use exact
field coordinates and recursion level.

Wolfram obtains the coordinates of \(A(\theta)\) from \code{AlgebraicNumber},
in ascending power order. The SymPy field representation uses
descending coefficient lists for Horner evaluation. Mixing these conventions
would silently corrupt automorphism actions. The accompanying implementations
make the order explicit.

\subsection{Exactness of the chosen embedding in Python}

The final Python equality check first computes monic minimal polynomials.
Different polynomials imply different values. Equal degree-one polynomials
identify the same rational value. For higher degree, equality of minimal
polynomials alone only proves conjugacy. The implementation additionally
calls \code{Poly.same\_root}, which uses a polynomial root-separation bound
and bounded-error evaluation to identify the selected root
\cite{sympyreference,sympysource}.

This is materially different from a fixed working-precision test. If a
separation bound is \(\delta>0\), approximations with certified errors small
relative to \(\delta\) can distinguish conjugates even when their first many
decimal digits agree. The regression suite includes
\[
        1+\sqrt2\,10^{-30}\quad\hbox{and}\quad1-\sqrt2\,10^{-30}.
\]
Their minimal polynomials are equal; their selected values are not.

The default real PSLQ path of SymPy's \code{to\_number\_field} checks an
exact defining-polynomial relation but uses a high-precision match to select
which conjugate that relation represents. For the package's phase and target
embeddings, \code{embed\_exact} instead calls
\code{field\_isomorphism} with \code{fast=False}. That factorization path selects
linear factors with the same-root test. The final public check is independent
of this coordinate construction. This distinction was found by inspecting
the installed SymPy source, not by assuming that every high-level call with
exact inputs automatically provides an independent branch certificate.

The package still trusts the CAS's minimal-polynomial and bounded-evaluation
implementations. Its \code{Verified} flag means a checked CAS computation,
not an exported proof object that a tiny standalone verifier or Lean kernel
can validate. Exporting full isolating-region and field-arithmetic
certificates would be a useful further project.

\Needspace{19\baselineskip}
\subsection{Failure taxonomy}

\begin{longtable}{@{}p{0.21\textwidth}p{0.73\textwidth}@{}}
\toprule Status & Meaning \\ \midrule\endhead
\code{Success} & The expression passed the radical grammar and selected-root check. \\
\code{NotSolvable} & An exact group computation proves the minimal polynomial's group nonsolvable. \\
\code{Unknown} & Fast mode found no certified structural conversion. No impossibility conclusion follows. \\
\code{ResourceLimit} & A time, degree, recursion or node budget was reached. No impossibility conclusion follows. \\
\code{BackendError} & A CAS operation failed, a required representation was unavailable, or an invariant did not verify. \\
\code{InvalidInput} & The Wolfram front end rejected the contract; the Python front end instead raises input-validation exceptions. \\
\bottomrule
\end{longtable}

In particular, \code{Unknown}, \code{ResourceLimit}, and \code{BackendError}
are never converted into \code{NotSolvable}. A finite unsuccessful search
cannot certify a negative mathematical statement.

\section{Complexity: finite does not mean small}

Let \(n=\deg f\), \(d=[L:\Q]\), \(D=[M:\Q]\), and \(h=|H|\).
Then
\[
 d\leq n!,\qquad D\leq d\varphi(\rad(d)),\qquad
 D=h\varphi(m).
\]
Even moderate root degree can therefore lead to a large primitive polynomial.
The degree cap is a guard on a mathematical quantity, not a complete memory
or bit-complexity bound. Rational coefficient heights can grow drastically
while the degree remains unchanged.

There is a useful bound on the \emph{Fourier expression recursion}, considered
separately. Ignoring memoization, a level with quotient prime \(p_i\) makes
at most \(p_i\) recursive calls. Hence the tree for one value has at most
\[
 1+p_r+p_rp_{r-1}+\cdots+\prod_{i=1}^r p_i
                   \leq 2h-1
\]
nodes, because every prime is at least two. At most \(h-1\) nontrivial
resolvent-root extractions occur in the full tree. This does not bound the
printed formula by \(O(h)\) characters: bottom coefficients and repeated
cyclotomic expressions can be large, and coefficient bit lengths are not
bounded by this count.

The reference implementation builds a full \(D\)-by-\(D\) automorphism
multiplication table, performs polynomial substitutions for field actions,
and checks fixed-field invariants repeatedly. With naive dense arithmetic,
these steps are expensive even before accounting for exact factorization,
primitive-element construction and embedding selection. No polynomial-time
claim in \(n\), input length, or output length is made.

The Python degree cap is checked after each successive field adjunction;
a single oversized adjunction can consume substantial resources before its
resulting degree is known. The Wolfram primitive-element call constructs
several generators together, so its wall-clock limit is particularly useful.
Python has no solver-internal wall-clock option. The supplied test runner
uses subprocess timeouts, and a production service should similarly isolate
expensive computations in a worker process.

The implemented fast paths are therefore substantive, not decorative. A
Dickson polynomial of degree seven or a reciprocal sextic can be solved
without constructing their full splitting fields. The general construction
is a fallback and an executable specification for more efficient future
backends.

\section{Using the supplied implementations}

\subsection{Wolfram Language}

Load the package directly from its extracted folder; a paclet descriptor is
also included. Direct loading avoids depending on paclet installation.
\begin{lstlisting}
Get["/path/to/RadicalRoots/Kernel/RadicalRoots.wl"];
r = Root[-1 - #^2 - #^3 + #^4 + #^6 &, 2];
report = RadicalRootReport[r, "Method" -> "Fast"];
report["Status"]
report["Expression"]
RadicalEqualQ[report["Expression"], r]
\end{lstlisting}
\code{RadicalRoot[r]} returns the expression on success and a \code{Failure}
object otherwise.

\code{RadicalRootReport[r]} returns an
association with status, method, verification flag and details. The default
method is \code{Automatic}; alternatives are \code{"Fast"} and
\code{"Galois"}. The public syntax and equality predicates are
\code{RadicalExpressionQ} and \code{RadicalEqualQ}.

The default budgets are 120 seconds, field degree 64, 10,000 Fourier nodes
and fast-recursion depth 30. Explicit extension hints are supported:
\begin{lstlisting}
RadicalRootReport[r,
  "ExtensionHints" -> {Sqrt[2], 3^(1/3)}]
\end{lstlisting}
The list identifies a compositum, not a list of separately attempted fields.
The hint-based factorization must still produce a certified smaller factor
that the fast layer can solve. To request the general algorithm without the
implemented budgets:
\begin{lstlisting}
RadicalRootReport[r, "Method" -> "Galois",
  "TimeConstraint" -> Infinity,
  "MaxFieldDegree" -> Infinity,
  "MaxNodes" -> Infinity,
  "MaxDepth" -> Infinity]
\end{lstlisting}
This can be extremely expensive. It does not disable backend limitations,
system memory limits or every possible internal CAS recursion limit.
Numbers already returned as radicals by initial normalization may bypass
the general construction even when that method is requested.

\subsection{Python and SymPy}

Install the pinned dependency and run the examples from the package root:
\begin{lstlisting}
python -m pip install -r python/requirements.txt
python examples/original_examples.py
python examples/embedded_model.py
python tests/run_tests.py --timeout 40
\end{lstlisting}
For direct use, put the package's \code{python} directory on \code{sys.path}
or \code{PYTHONPATH}:
\begin{lstlisting}
import sympy as s
from radicalroots import to_radicals
x = s.Symbol("x")
r = to_radicals(x**6+x**4-x**3-x**2-1, 1,
                method="fast")
assert r.status == "Success" and r.verified
print(r.expression)
\end{lstlisting}
The methods are \code{"auto"}, \code{"fast"}, and \code{"galois"}.
\code{None} disables each of \code{max\_field\_degree},
\code{max\_nodes}, and \code{max\_depth}. There is no internal time option.
The result is a \code{RadicalResult} dataclass with \code{status},
\code{expression}, \code{method}, \code{verified}, and \code{details}.

The standalone automatic model does not require a special Galois-group
function for arbitrary degree. It computes its own embedded automorphisms
from the splitting field. SymPy's group routine, documented for irreducible
rational polynomials through degree six, is only an optional early negative
decision accelerator \cite{sympynumberfields}. Thus this limit is not an
artificial degree-six restriction of the mathematical backend.

\subsection{Supplying a known embedded field}

A field prepared elsewhere can avoid expensive automatic discovery:
\begin{lstlisting}
from galois_core import (GaloisModel, radicals_from_model,
                        embed_exact)
K = s.QQ.algebraic_field(s.sqrt(2), s.sqrt(3))
model = GaloisModel.from_field(K, conductor=2)
a = embed_exact(K, s.sqrt(2)-s.sqrt(3))
e, details = radicals_from_model(model, a)
\end{lstlisting}
This interface requires an actual embedded Galois field containing the
specified principal root of unity. The conductor must include every prime
dividing the relative group order. An abstract group or a list of approximate
roots does not satisfy the contract.

\section{Validation record and limitations}
\label{sec:validation}

The final isolated regression run completed \textbf{36 tests: 36 passed, zero failed, and zero timed out}, with a 40-second limit per test.
The machine-readable \path{test-results/python-tests.json} records the actual
Python version, SymPy version, platform, subprocess timeout, elapsed time and
output for each regression. The tests cover the following distinct issues:
\begin{itemize}
\item Every conjugate of the motivating sextic and quintic, plus low-degree,
shifted-binomial, degree-seven Dickson and composition identities.
\item Full automatic Fourier conversion for quadratic, cubic and a
nonabelian degree-eight splitting-field example; supplied real and complex
biquadratic models; cyclotomic-base-only descent.
\item Nonzero phase corrections, zero resolvents, near-coincident conjugates,
strict output syntax, invalid input, invalid field models, resource limits,
and exact nonsolvability of an \(S_5\) example.
\end{itemize}
The finite-group component is independently exercised on \(C_5,S_3,S_4\),
the dihedral group of order eight, and the perfect group \(A_5\).
The \(S_5\) negative test uses \(X^5-X-1\) and checks the exact reported
order 120. The reducible-input test additionally selects a rational root
from a polynomial containing that nonsolvable factor.

The final test-result file is authoritative for pass counts and timings;
\code{README.md} summarizes it. Tests are not a proof of universal program
correctness. In particular, the suite does not establish practical automatic
splitting-field performance on arbitrary solvable quintics, septics or the
large polynomials in the notebook. Early exploratory adjunctions of multiple
\code{CRootOf} generators were slow enough to motivate the low-degree
internal-radical optimization. No claim is made to have run the complete
notebook or to have benchmarked a competitive high-degree Galois solver.

The native Wolfram kernel was unavailable in the authoring environment.
The connected Wolfram service returned HTTP 404 before evaluation on both
attempts. Therefore the supplied \code{.wlt} tests are \emph{not} reported as
passed. A static delimiter/comment/string check is useful for packaging,
but is not a substitute for executing the package. The Python results do
not certify the Wolfram translation.

There are also intentional scope limitations. The solver does not denest
or optimize a returned tower globally. Real inputs may produce complex
intermediate radicals. It does not provide an all-real-radical guarantee,
which is a stronger condition. Input recognition from numerical constants,
parameter stratification and proofs about inverse special functions remain
outside its contract. Failure reports preserve these distinctions rather
than silently returning an unchecked expression.

\section{Related work and practical next steps}

The theoretical ingredients are classical: solvability of finite groups,
Galois correspondence, cyclotomic base change and cyclic-extension
resolvents. A recent explicit algorithmic treatment by Song Li also uses
Lagrange/Fourier-resolvent ideas \cite{li}. The present package is an
independently written reference implementation organized around actual
embedded-field actions, a concrete compositum that fixes the needed roots
of unity, and explicit principal-branch selection. It does not adopt
unverified runtime claims for a different representation or implementation.

Several optimizations can preserve the exact contract. A stronger backend
could construct relative fields rather than a large absolute primitive
field, use group stabilizers and subfield algorithms to find smaller towers,
or compute only the field actions needed for one target. Recognized
cyclotomic periods and solvable quintic normal forms could be additional
fast paths. Such accelerators should propose exact structures, then pass the
same invariant and branch checks; none should replace a failed search with
a nonsolvability claim.

A further useful separation is an independently checkable certificate:
primitive polynomials and isolating regions, rational matrices for
automorphisms, subgroup inclusions, resolvent powers, and branch-isolating
certificates. Those data could allow an external verifier to check an answer
without repeating field discovery. This release records decisions and
verifies them in the CAS, but it does not yet export that full proof format.

The central conclusion is practical as well as theoretical. The original
examples do not require blind searches over complicated extension guesses;
their structures give short radicals directly. For the remaining solvable
cases, a complete constructive route is available. Its main obstacle in a
general-purpose implementation is the cost and robustness of exact field
computation, not the absence of a mathematical algorithm.

\appendix
\section{Package map and reproducibility}

\begin{longtable}{@{}p{0.42\textwidth}p{0.52\textwidth}@{}}
\toprule Path & Purpose \\ \midrule\endhead
\code{Kernel/RadicalRoots.wl} & Wolfram package and both conversion layers. \\
\code{Kernel/init.m}, \code{PacletInfo.wl} & Loader and optional paclet metadata. \\
\code{python/radicalroots.py} & Public Python interface, structural proposals and selected-root checks. \\
\code{python/galois\_core.py} & Exact field models, automorphisms, prime series and Fourier descent. \\
\code{python/requirements.txt} & Pinned SymPy dependency. \\
\code{tests/test\_python.py} & Executable regressions. \\
\code{tests/run\_tests.py} & Isolated runner with actual JSON results. \\
\code{tests/RadicalRoots.wlt} & Native Wolfram regression specification. \\
\code{tests/run\_wolfram.wls} & Native test runner. \\
\code{examples/} & Original examples and supplied-model demonstration. \\
\code{test-results/} & Actual Python results and explicit Wolfram execution status. \\
\code{source-review/} & Notebook inspection, transcription, provenance and hashes. \\
\code{docs/article.tex}, \code{docs/article.pdf} & This article. \\
\bottomrule
\end{longtable}

Compile the article from its directory with two runs of
\code{pdflatex -interaction=nonstopmode -halt-on-error article.tex}.
The bibliography is embedded; no BibTeX invocation is required. Runtime
measurements include test-process startup and should not be read as isolated
solver benchmarks. The new implementation is MIT-licensed; the original
notebook-derived material retains its separate attribution.

\begingroup\raggedright
\begin{thebibliography}{9}
\bibitem{question}
Vladimir Reshetnikov, \emph{Why is ToRadicals not able to handle all cases?
Is there a workaround?}, Mathematica Stack Exchange, question 34011.
\url{https://mathematica.stackexchange.com/questions/34011/why-is-toradicals-not-able-to-handle-all-cases-is-there-a-workaround}.

\bibitem{toradicals}
Wolfram Research, \emph{ToRadicals}, Wolfram Language documentation.
\url{https://reference.wolfram.com/language/ref/ToRadicals.html}.

\bibitem{tonumberfield}
Wolfram Research, \emph{ToNumberField} and \emph{RootReduce}, Wolfram Language documentation.
\url{https://reference.wolfram.com/language/ref/ToNumberField.html};
\url{https://reference.wolfram.com/language/ref/RootReduce.html}.

\bibitem{milne}
J. S. Milne, \emph{Fields and Galois Theory}, version 5.10, 2022,
especially Chapter 5 and Theorem 5.34.
\url{https://www.jmilne.org/math/CourseNotes/FT.pdf}.

\bibitem{sympynumberfields}
The SymPy development team, \emph{Number Fields}, SymPy 1.14.0 documentation.
\url{https://docs.sympy.org/latest/modules/polys/numberfields.html}.

\bibitem{sympyreference}
The SymPy development team, \emph{Polynomials Manipulation Module Reference},
entry \code{Poly.same\_root}, SymPy 1.14.0 documentation.
\url{https://docs.sympy.org/latest/modules/polys/reference.html}.

\bibitem{sympysource}
The SymPy development team, SymPy 1.14.0 source code,
\code{sympy/polys/numberfields/subfield.py} and
\code{sympy/polys/polytools.py}. The installed source was inspected for
\code{field\_isomorphism\_pslq}, \code{field\_isomorphism\_factor},
and \code{Poly.same\_root}.
\url{https://github.com/sympy/sympy/tree/1.14.0/sympy/polys}.

\bibitem{li}
Song Li, \emph{An Algorithm for Solving Solvable Polynomial Equations of Arbitrary Degree by Radicals}, arXiv:2203.11602v2, 2022.
\url{https://arxiv.org/abs/2203.11602}.

\end{thebibliography}
\endgroup
\small Online references consulted 8 September 2026. The supplied notebook is
identified separately by its archive and notebook SHA-256 hashes in
\code{source-review/README.md}.
\end{document}
